% 合成示例（synthetic example），非正式案例材料：plain article 原稿。
% 故意埋入两处宏包问题，供 pkgconflict.py 检出：
%   1) geometry 与目标模板类自带 \RequirePackage[margin=2.5cm]{geometry} 重复且选项冲突；
%   2) hyperref 未放在稿件宏包清单的最后。
\documentclass{article}
\usepackage[margin=3cm]{geometry}
\usepackage{hyperref}
\usepackage{booktabs}
\usepackage{graphicx}

\title{Adaptive Mesh Refinement for Shallow Water Simulations}
\author{Jane Doe \and John Smith}
\date{\today}

\begin{document}
\maketitle

\begin{abstract}
We present an adaptive mesh refinement strategy for shallow water
simulations. The method refines cells near steep gradients and keeps the
coastal zones coarse. Numerical tests show a factor of three speedup
without loss of accuracy.
\end{abstract}

\section{Introduction}
Shallow water equations describe free-surface flows in rivers and coastal
regions \cite{leveque2002}. Uniform fine grids are expensive, so adaptive
strategies concentrate resolution where the solution changes rapidly
\cite{berger1984}.

\section{Numerical Results}
Table~\ref{tab:speedup} summarizes the measured speedup on two benchmark
cases. The adaptive grid reduces the cell count while preserving the
peak wave height.

\begin{table}[htbp]
  \centering
  \caption{Speedup of the adaptive grid on two benchmark cases.}
  \label{tab:speedup}
  \begin{tabular}{lcc}
    \toprule
    Case & Cells & Speedup \\
    \midrule
    River bend & 12000 & 2.8 \\
    Coastal bay & 45000 & 3.1 \\
    \bottomrule
  \end{tabular}
\end{table}

The error in peak wave height stays below one percent in both cases, so
the refinement criterion is accepted for production runs.

\begin{thebibliography}{9}
\bibitem{leveque2002}
R.~J. LeVeque, \emph{Finite Volume Methods for Hyperbolic Problems},
Cambridge University Press, 2002.

\bibitem{berger1984}
M.~J. Berger and J.~Oliger, Adaptive mesh refinement for hyperbolic
partial differential equations, \emph{Journal of Computational Physics}
53 (1984), 484--512.
\end{thebibliography}

\end{document}
